\chapter{Raft as a Pluggable Mako Backend}

Chapter 2 described the execution and replication shape that the original
Paxos path gave to Mako. This chapter explains how Raft is placed behind that
same replication boundary. The goal is not to retell the Raft protocol. The
goal is to explain the systems integration: how Mako continues to submit
transaction log records through the same helper interface, how those records
enter Raft, and why the implementation evaluates both Multi-Raft and
Single-Raft topologies.

\section{Integration Goals}

The first goal is API compatibility. Mako's transaction execution layer should
not need to know whether Paxos or Raft orders a transaction log record. Worker
threads still call the same replication helper functions to set up replication,
register callbacks, submit log records, wait for submission, and shut down the
backend. The Raft integration therefore sits behind the existing helper
boundary rather than replacing Mako's transaction execution engine.

The second goal is experimental comparability. A Paxos/Raft comparison is only
useful if the outer experiment is held constant. The Raft backend runs under
the same benchmark driver, worker-thread counts, TPC-C workload, CPU
measurement path, and result extraction scripts as the Paxos backend. The
replication backend changes, but the workload generator and the thesis metrics
remain the same.

The third goal is pipeline preservation. Raft is not used as an opaque log
placed after Mako. It participates in the same execution path: Mako workers
produce transaction log records, the replication helper submits them to the
configured backend, the consensus layer commits them, and committed entries
are delivered back through callbacks. The integration is successful only if
this path preserves both leader-side submission and follower-side replay
progress.

\section{Runtime Replication Boundary}

The codebase uses a dispatch layer to keep Mako protocol-agnostic. The shared
replication helper stores a runtime replication type and forwards each helper
operation to either the Paxos implementation or the Raft implementation. From
Mako's point of view, the important operations still have the same outer
meaning: initialize the replication backend, register the leader-side and
follower-side callbacks, submit a transaction log record, wait for outstanding
submissions when needed, and shut the backend down.

This boundary is important because it keeps the comparison honest. The
transaction layer does not special-case Raft. The benchmark can select the
replication backend at startup, and the helper dispatches every later call to
the selected implementation. In the Raft case, the selected implementation
creates Raft-side workers that submit transaction log records into Raft. In
the Paxos case, the selected implementation uses the original Paxos workers.

The same boundary also limits what Chapter 3 needs to explain. The body of the
thesis does not need a file-by-file inventory of every helper function. The
relevant architectural point is simpler: Mako submits and receives committed
work through one stable replication interface, while the backend behind that
interface can be Paxos, Multi-Raft, or Single-Raft.

\section{Topology Choices: Multi-Raft and Single-Raft}

Using the worker-lane terminology from Chapter 2, the Raft integration has two
topology choices. Multi-Raft keeps the original Paxos-shaped structure:
separate partition lanes are backed by separate Raft groups. Single-Raft
consolidates those lanes into one ordered Raft log and carries partition
identity inside each replicated command.

\figref{fig:replication-topologies} shows the difference. The left panel is
the original Multi-Paxos shape. The middle panel is Multi-Raft: each partition
lane maps to its own Raft group, so the replication topology preserves the
same broad parallelism as the original backend. The right panel is
Single-Raft: partition lanes submit into one consolidated Raft group, with one
local replica of that group running in each benchmark process. After commit,
the partition-lane identifier inside the command tells Mako which lane should
receive the committed record.

\begin{figure}[!htbp]
\centering
\includegraphics[width=0.96\linewidth]{figures/manual/fig02_replication_topologies.pdf}
\caption[Replication topologies]{Replication topologies evaluated in this thesis. Multi-Raft preserves the original per-partition parallelism with one Raft group per partition lane. Single-Raft consolidates partition commands into one Raft log, with one replica of the consolidated group running in each process.}
\label{fig:replication-topologies}
\end{figure}

Multi-Raft is the conservative design. It asks whether Raft can remain in
Paxos's performance class when it is given a topology that matches the
original system. Each local partition lane has a corresponding
Raft-side worker and local Raft replica. Work from different lanes does not
share one Raft log, so the design preserves the same kind of replication
parallelism that made the original Multi-Paxos path a natural fit for Mako.

Single-Raft asks a different question. One consolidated Raft group spans the
replica processes, and each process hosts one local replica for that group.
All partition commands from a process pass through that shared log. This
reduces the amount of Raft machinery, but it changes the pressure point.
Instead of many independently ordered partition logs, the process now has one
ordered path. Because many Mako lanes share that path, each committed record
must carry enough information to be delivered back to the lane that produced
it. That makes callback routing and follower replay engineering central to
the Single-Raft design.

\section{Raft Worker Setup}

The Raft-side worker is the bridge between Mako's transaction-log pipeline and
the Raft protocol implementation. It occupies the same architectural position
as the original Paxos worker: Mako hands it serialized transaction log records,
the worker submits those records to consensus, and committed records come back
through an application callback.

Raft worker setup has four conceptual steps. First, the worker creates the
Raft protocol stack for its local replica. The replica is configured with the
local site information and the partition lane it serves, so it can join the
right group. Second, the worker starts the RPC service used by Raft. This
service receives election messages and replicated-log messages from other
replicas. Third, the worker creates the communication path used to send Raft
messages to peers. Fourth, the worker starts supporting control paths, such as
heartbeat handling, and starts the background submission thread that drains
pending Mako log records.

Multi-Raft and Single-Raft differ in how many workers are created. In
Multi-Raft, setup creates one worker for each local partition lane. In
Single-Raft, setup creates one worker from the first local site configuration
and stores the other site bindings so the process can preserve the network
shape expected by the rest of the configuration. The Single-Raft path creates
stub RPC servers for the extra partition ports; these stubs preserve expected
connections for the existing configuration, but they do not own a replicated
log, vote as independent Raft replicas, or apply committed commands. The
single real local Raft replica owns the consolidated Raft log.

Both designs also configure preferred leadership during setup. Chapter 4 gives
the full design and safety argument. For this chapter, the important point is
that preferred leader configuration is part of making Raft fit Mako's
leader-oriented submission path, and it applies to both Multi-Raft and
Single-Raft.

\section{Log Submission Path}

The log submission path answers a simple question: after a Mako worker finishes
a transaction, how does the resulting transaction log record enter Raft? At
this point the transaction has already been executed by Mako. The Raft backend
does not rerun the transaction. Its job is to take the record that Mako
produced, put that record into a replicated Raft log, and later tell Mako when
the record has committed.

A concrete example helps. Suppose the worker for lane 3 finishes a transaction
and produces a transaction log record. In Multi-Raft, lane 3 has its own Raft
group, so the record is submitted to the lane-3 Raft group. The group itself
remembers where the record came from: any committed record from that group
belongs to lane 3. In Single-Raft, there is only one consolidated Raft group
for all lanes, so the backend submits the same record to the shared Raft log.
Because the shared log contains records from many lanes, the record must carry
an explicit lane label saying that it came from lane 3.

Before the backend accepts the record, it checks leadership. Raft only accepts
new client commands at the leader of the group that will order the record. In
Multi-Raft, the relevant leader is the leader of the lane-3 Raft group. In
Single-Raft, the relevant leader is the leader of the one consolidated group.
If the local process is not that leader, the backend refuses the local
submission instead of pretending that a follower can accept leader-side work.
For the benchmark, that refusal means the operation is not counted as committed
work. Chapter 4 explains how preferred leader election makes the intended
leader-side process hold leadership in the normal case.

Once the leader check passes, the backend makes its own copy of the transaction
log bytes and places that copy in an internal submission queue. This copy is a
practical ownership step: the Mako worker can return and reuse its temporary
buffer, while Raft keeps a stable copy to replicate. A background submission
thread drains the queue and feeds records into the local Raft replica. This
keeps the worker-facing path short and prevents Mako worker threads from doing
all of the Raft append work inline.

The last step before Raft sees the record is packaging. Raft replicates
commands, so the backend wraps the transaction log bytes as a Raft command. In
Multi-Raft, the command only needs the transaction log bytes because the Raft
group already identifies the lane. In Single-Raft, the command contains both
the transaction log bytes and the partition-lane identifier. That extra label
is what makes the later callback delivery in Section~\ref{sec:commit-callback}
possible: when the shared log commits the command, the backend can still route
the committed record back to lane 3.

\section{Commit and Callback Delivery}
\label{sec:commit-callback}

Commit and callback delivery are easier to understand if they are separated
into two questions. First, when does Raft decide that a submitted record is
committed? Second, after that decision, which Mako path receives the committed
record?

For the first question, consider a Mako worker that produced a transaction log
record for lane 3. The local Raft leader receives the wrapped command for that
record and appends it to its own Raft log. At this moment the record has a
position in the leader's log, but it is still only proposed work. Mako should
not replay it yet, because the record has not been accepted by enough replicas
to survive a leader change.

The leader then sends the log entry to follower replicas. Each follower checks
that the new entry is consistent with the log prefix it already has. Once a
majority of replicas has accepted entries far enough for Raft's commit rule,
the leader can advance its commit point. In Raft, committing an entry from the
current term can also make earlier entries committed indirectly. For this
thesis, the important idea is simple: the record becomes committed only after
Raft has enough replicated evidence that the record belongs in the agreed log.

The second question is routing the committed record back into Mako. Raft has
chosen the order, but the Raft backend still has to deliver the committed
transaction log record to the correct part of Mako's database pipeline. This
delivery happens through a callback registered by Mako. On the leader replica,
the callback is used for leader-side completion: it tells the leader path that
the submitted record has become committed. On follower replicas, the callback
is used for follower-side replay: it passes the committed record to the replay
path so the follower's local database state can catch up.

The routing rule differs between Multi-Raft and Single-Raft. In Multi-Raft,
lane 3 has its own Raft group, so a committed record from that group already
identifies lane 3. In Single-Raft, all lanes share one consolidated Raft log,
so the group alone is not enough information. The committed command must carry
the partition-lane identifier that was attached during submission. When the
Single-Raft apply path receives the committed command, it reads that identifier
and delivers the record to the matching Mako callback for lane 3.

This is the boundary of Chapter 3: Raft commits an ordered record and hands it
back to Mako through the correct callback. What happens inside the callback is
the next design question. It may be lightweight staging, or it may trigger
heavier follower replay work. Chapter 5 explains why Single-Raft needed
ReplayPool to keep that heavier replay work out of the narrow Raft apply path.

\section{What Chapter 3 Establishes}

The Raft integration is a pluggable backend, not a separate benchmark path.
Mako still submits transaction log records through the shared replication
interface, and the selected backend determines whether those records are
ordered by Paxos, by per-partition Raft groups, or by a consolidated Raft log.

Multi-Raft is the direct Paxos-shaped replacement. It keeps one Raft group per
partition lane and therefore preserves the original replication parallelism.
Single-Raft is the consolidation experiment. It uses one process-wide worker
and one consolidated Raft log to order all partition commands, so it must route
committed entries back to their original partition lanes and must handle
follower replay carefully.

These two designs set up the next two chapters. Chapter 4 explains why both
Raft designs need preferred leader election to match Mako's leader-side
submission path. Chapter 5 explains why Single-Raft also needs ReplayPool to
preserve follower replay progress after consolidation.
